% !TEX encoding = UTF-8
% !TEX program = xelatex
% !TEX spellcheck = en_US

% English counterpart of contrib/latex-parity/main.typ. This
% entrypoint is compiled from
% the pinned ThuThesis checkout and deliberately uses its support files.
\providecommand\CompareOutput{electronic}
\documentclass{thuthesis}
\input{thusetup}

\thusetup{
  output = \CompareOutput,
  language = english,
  title = {Tsinghua University Typst Thesis Template\\Sample Document v0.1.0},
  title* = {An Introduction to the Typst Thesis Template of Tsinghua
  University},
  degree-category = {Master of Science},
  degree-category* = {Master of Science},
  department = {Department of Computer Science and Technology},
  discipline = {Computer Science and Technology},
  discipline* = {Computer Science and Technology},
  author = {Xue Ruini},
  author* = {Xue Ruini},
  supervisor = {Zheng Weimin, Professor},
  supervisor* = {Professor Zheng Weimin},
  associate-supervisor = {Chen Wenguang, Professor},
  associate-supervisor* = {Professor Chen Wenguang},
  date = {2026-08-08},
  clc = {TP3},
  udc = {004},
  id = {2026-01},
  discipline-level-1 = {Computer Science and Technology},
  discipline-level-2 = {Computer Systems Architecture},
  start-date = {2024-08-01},
  end-date = {2026-07-31},
}

\begin{document}

\maketitle

\begin{committee}[name={Thesis Supervision Committee, Reviewers, and
  Defense Committee}]
  \newcolumntype{C}[1]{@{}>{\centering\arraybackslash}p{#1}}
  \section*{Thesis Supervision Committee}
  \begin{center}
    \begin{tabular}{C{3cm}C{4cm}C{8cm}@{}}
      Li XX    & Professor           & Tsinghua University \\
      Wang XX  & Associate Professor & Tsinghua University \\
      Zhang XX & Assistant Professor & Tsinghua University \\
    \end{tabular}
  \end{center}
  \section*{External Reviewers}
  \begin{center}Anonymous review
  \end{center}
  \section*{Defense Committee}
  \begin{center}
    \begin{tabular}{C{2.75cm}C{2.98cm}C{4.63cm}C{4.63cm}@{}}
      Chair     & Zhao XX & Professor           & Tsinghua University \\
      Member    & Liu XX  & Professor           & Tsinghua University \\
      Member    & Yang XX & Research Professor  & Institute XXXXXXX,
      Chinese Academy of Sciences \\
      Secretary & Wu XX   & Assistant Research Professor & Tsinghua
      University \\
    \end{tabular}
  \end{center}
\end{committee}

\copyrightpage

\frontmatter
\input{data/abstract}
\tableofcontents
\listoffigures
\listoftables

\begin{denotation}[3cm]
\item[PI] Polyimide
\item[MPI] Polyimide model compound, N-phenylphthalimide
\item[PBI] Polybenzimidazole
\item[MPBI] Polybenzimidazole model compound, N-phenylbenzimidazole
\item[PY] Polypyrrolone
\item[PMDA-BDA] Polypyrrolone film synthesized from PMDA and BDA
\item[MPY] Polypyrrolone model compound
\item[As-PPT] Asymmetric polyphenyl triazine
\item[MAsPPT] Asymmetric polyphenyl triazine monomer model compound
\item[DMAsPPT] Asymmetric polyphenyl triazine dimer model compound
\item[S-PPT] Symmetric polyphenyl triazine
\item[MSPPT] Symmetric polyphenyl triazine model compound
\item[PPQ] Polyphenylquinoxaline
\item[MPPQ] Polyphenylquinoxaline model compound
\item[HMPI] Protonated product of the polyimide model compound
\item[HMPY] Protonated product of the polypyrrolone model compound
\item[HMPBI] Protonated product of the polybenzimidazole model compound
\item[HMAsPPT] Protonated product of the asymmetric polyphenyl
  triazine model compound
\item[HMSPPT] Protonated product of the symmetric polyphenyl triazine
  model compound
\item[HMPPQ] Protonated product of the polyphenylquinoxaline model compound
\item[PDT] Thermal decomposition temperature
\item[HPLC] High-performance liquid chromatography
\item[HPCE] High-performance capillary electrophoresis
\item[LC-MS] Liquid chromatography--mass spectrometry
\item[TIC] Total ion content
\item[\textit{ab initio}] First-principles quantum-chemistry calculation method
\item[DFT] Density functional theory
\item[$E_a$] Activation energy of a chemical reaction
\item[ZPE] Zero-point vibrational energy
\item[PES] Potential energy surface
\item[TS] Transition state
\item[TST] Transition-state theory
\item[$\increment G^\neq$] Activation free energy
\item[$\kappa$] Transmission coefficient
\item[IRC] Intrinsic reaction coordinate
\item[$\nu_i$] Imaginary frequency
\item[ONIOM] Our Own N-layered Integrated Molecular Orbital and
  Molecular Mechanics
\item[SCF] Self-consistent field
\item[SCRF] Self-consistent reaction field
\end{denotation}

\mainmatter

\chapter{Writing the Main Body of a Thesis}

In addition to satisfying formal presentation requirements, a graduate thesis
should follow several common principles in its content.

A graduate thesis should normally address one coherent theme rather than combine
several unrelated pieces of work. It should investigate a specific problem in a
discipline thoroughly and systematically and reach conclusions of value. The
research theme should not be excessively broad. For example, ``A Study of the
Reform of Chinese State-owned Enterprises'' is too broad because such reform
involves too many issues to examine fully in a single graduate thesis.

\section{Language and Expression}

Except for international graduate students, degree theses must be written in
Chinese. They should use standardized Chinese characters and simplified Chinese,
apart from ancient characters required by research on classical Chinese and
foreign-language works cited in the references.

International graduate students normally write their theses in Chinese or
English and follow the same formatting requirements. The thesis must have a
Chinese cover.

A graduate thesis is an academic work. Its exposition should therefore be
rigorous and concise, with clear emphasis. Common disciplinary knowledge should
be stated briefly or omitted. The argument should be sound, the data reliable,
the explanation thorough, the reasoning rigorous, the wording economical, and
the organization clear. Literary, emotive, or otherwise nonacademic language
should be avoided.

Any uncommon new name, term, or concept introduced in the thesis should be
explained immediately.

\section{Writing the Thesis Title}

The title should reflect the principal content of the thesis concisely,
precisely, and specifically. It is an accurate description of the object of
study and should generally indicate the conclusion to some extent. A title
should tell readers not only what problem the thesis investigates but also what
the study concludes. For example, ``Between Fact and Fiction: The Views of Mel,
Carpenter, and Wolfe on Journalism'' is more professional and precise than ``A
Study of Three American Writers' Views on Journalism.''

\section{Writing the Abstract}

The abstract is a highly condensed account of the thesis and should be
independent and self-contained: a short but complete article. It should give the
reader an overall understanding of the research methods and conclusions.
Accordingly, it should describe the problem and objective, briefly introduce the
methods and research process, and distill the conclusions, with emphasis on the
results and conclusions.

The abstract should not be written as an outline of the full text. In
particular,
statements such as ``Chapter 1~\ldots; Chapter 2~\ldots; \ldots'' should be
avoided.

\section{Writing the Introduction}

The introduction of a thesis generally contains five parts: (1) statement of the
problem; (2) background and significance; (3) literature review; (4) research
methods; and (5) organization of the thesis.
\begin{itemize}
  \item \textbf{Statement of the problem:} Clearly explain what
    problem is being studied.\footnote{Choose a topic with a genuine
    problem to solve; do not select a non-problem as the object of research.}
  \item \textbf{Background and significance:} Explain why the topic
    was selected, including its contribution to the discipline and
    its theoretical and practical significance.
  \item \textbf{Literature review:} Synthesize and critically review
    work within the topic area. Identify the current state of
    research, unresolved problems, and the exploratory contribution
    of the present study.
  \item \textbf{Research methods:} Explain the academic research
    methods used in the thesis.
  \item \textbf{Organization of the thesis:} Introduce the structure
    of the document.
\end{itemize}

\section{Writing the Main Text}

This part presents the author's research and should distinguish it clearly from
the work of others. Material already discussed in the literature review need not
be repeated. The chapters should be organically connected and follow a logical
sequence.

\section{Writing the Conclusions}

The conclusions distill and summarize the main results and arguments.
They should
be concise, accurate, and complete, enabling readers to understand the
significance, purpose, and work of the thesis. They are the final overall
conclusions, not a simple repetition of chapter summaries. They should state the
central arguments, explain the author's original work and the
position, role, and
significance of the results in the field, identify limitations, and offer
suggestions for future work. The author's results must be distinguished strictly
from those of the supervisor and other researchers.

Assessment of the work should remain factual. Terms such as
``first,'' ``leading,''
or ``fills a gap'' should be avoided unless supported by sufficient evidence.

\chapter{Examples of Figures and Tables}

\section{Figures}

Images are normally inserted in a \texttt{figure} environment with
\verb|\includegraphics|, as in Figure~\ref{fig:example-en}. Vector graphics
should preferably use PDF, photographs should use JPG, and other raster graphics
should use lossless PNG.

\begin{figure}
  \centering
  
\includegraphics[width=0.5\linewidth]{example-image-a.pdf}
  \caption*{Note: Keep the title separate from explanatory text,
  which belongs below the figure or in the main text.}
  \caption{Example image}
  \label{fig:example-en}
\end{figure}

Notes to a figure or table should be identified by lowercase English letters in
sequence and placed below it. If a figure contains two or more subfigures, label
them (a), (b), (c), and so on, and provide a title for each one.

\begin{figure}
  \centering
  \subcaptionbox{Subfigure
  A}{
\includegraphics[width=0.35\linewidth]{example-image-a.pdf}}
  \subcaptionbox{Subfigure
  B}{\includegraphics[width=0.35\linewidth]{example-image-b.pdf}}
  \caption{Example with multiple subfigures}
\end{figure}

\section{Tables}

A table should be self-explanatory, and every quantity and unit should be
identified by its symbol. For clarity and readability, use three-line tables
such as Table~\ref{tab:three-line-en}.

\begin{table}
  \centering
  \caption{Example three-line table}
  \begin{tabular}{ll}
    \toprule
    File name & Description \\
    \midrule
    thuthesis.dtx   & Template source, including documentation and comments \\
    thuthesis.cls   & Template class file \\
    thuthesis-*.bst & BibTeX bibliography style files \\
    \bottomrule
  \end{tabular}
  \label{tab:three-line-en}
\end{table}

Notes in a table should be identified by lowercase English letters in sequence
and placed below the table.

\begin{table}
  \centering
  \begin{threeparttable}[c]
    \caption{Example table with notes}
    \begin{tabular}{ll}
      \toprule
      File name & Description \\
      \midrule
      thuthesis.dtx\tnote{a} & Template source, including
      documentation and comments \\
      thuthesis.cls\tnote{b} & Template class file \\
      thuthesis-*.bst & BibTeX bibliography style files \\
      \bottomrule
    \end{tabular}
    \begin{tablenotes}
    \item[a] This file can generate the template manual or extract a
      compact class file.
    \item[b] Regenerate the class file when updating the template.
    \end{tablenotes}
  \end{threeparttable}
\end{table}

When a table continues onto another page, mark it as a continued table and
repeat the header on every page.

\begin{longtable}{cccc}
  \caption{Title of a multipage table} \\
  \toprule
  Heading 1 & Heading 2 & Heading 3 & Heading 4 \\
  \midrule
  \endfirsthead
  \caption*{Continued table~\thetable} \\
  \toprule
  Heading 1 & Heading 2 & Heading 3 & Heading 4 \\
  \midrule
  \endhead
  \bottomrule
  \endfoot
  Row 1 & & & \\ Row 2 & & & \\ Row 3 & & & \\ Row 4 & & & \\
  Row 5 & & & \\ Row 6 & & & \\ Row 7 & & & \\ Row 8 & & & \\
  Row 9 & & & \\ Row 10 & & & \\
\end{longtable}

\section{Algorithms}

The template does not prescribe a pseudocode package. The following example uses
the \texttt{algorithms} package.

\renewcommand{\algorithmicrequire}{\textbf{Input:}\unskip}
\renewcommand{\algorithmicensure}{\textbf{Output:}\unskip}
\begin{algorithm}
  \caption{Compute $y=x^n$}
  \small
  \begin{algorithmic}
    \REQUIRE $n \geq 0$
    \ENSURE $y=x^n$
    \STATE $y \leftarrow 1$
    \STATE $X \leftarrow x$
    \STATE $N \leftarrow n$
    \WHILE{$N \neq 0$}
    \IF{$N$ is even}
    \STATE $X \leftarrow X \times X$
    \STATE $N \leftarrow N/2$
    \ELSE
    \STATE $y \leftarrow y \times X$
    \STATE $N \leftarrow N-1$
    \ENDIF
    \ENDWHILE
  \end{algorithmic}
\end{algorithm}

\chapter{Mathematical Symbols and Equations}

\section{Mathematical Symbols}

Mathematical symbols in a Chinese thesis normally follow GB/T 3102.11--1993,
``Mathematical Signs and Symbols for Use in the Physical Sciences and
Technology,''\footnote{The former mandatory standard became a
recommended standard on March 23, 2017.}
which adopted ISO 31-11:1992\footnote{The current edition is ISO 80000-2:2019.}
with modifications. Its conventions differ in several respects from the
defaults used by the American Mathematical Society.

The author should apply the chosen standard consistently to mathematical
constants, special functions, differential symbols, vectors, and matrices. For
example,
\begin{equation*}
  \uppi=3.14\dots; \quad \symup{i}^2=-1; \quad
  \symup{e}=\lim_{n\to\infty}\left(1+\frac{1}{n}\right)^n.
\end{equation*}
Vectors, matrices, and tensors are normally bold italic, as in $\symbf{x}$ and
$\symbf{\Sigma}$; the natural logarithm is written $\ln x$.

English theses use the default \TeX{} mathematical style. Authors may select a
different style explicitly when required. Quantities and units should use
consistent formats, for example \SI{6.4e6}{m}, \SI{9}{\micro\meter},
\si{kg.m.s^{-1}}, and \SIrange{10}{20}{\degreeCelsius}.

\section{Mathematical Equations}

Display equations may use the \texttt{equation} environment. References should
be enclosed in parentheses, as in Equation~\eqref{eq:example-en}.
\begin{equation}
  \frac{1}{2\uppi\symup{i}}\int_\gamma f
  =\sum_{k=1}^m n(\gamma;a_k)\mathscr{R}(f;a_k).
  \label{eq:example-en}
\end{equation}
Multiline equations should align at the equals sign where possible:
\begin{align}
  a &= b+c+d+e \\
  &= f+g.
\end{align}

\section{Mathematical Theorems}

The template can configure theorem environments through \texttt{amsthm}.
\begin{theorem}[Lindeberg--L\'{e}vy central limit theorem]
  Let $X_1,X_2,\dots,X_n$ be independent and identically distributed random
  variables with expectation $\mu$ and finite variance $\sigma^2\ne0$. Define
  $\bar X_n=n^{-1}\sum_{i=1}^n X_i$. Then
  \begin{equation}
    \lim_{n\to\infty}P\left(\frac{\sqrt n(\bar X_n-\mu)}{\sigma}\le
    z\right)=\Phi(z),
  \end{equation}
  where $\Phi(z)$ is the cumulative distribution function of the standard normal
  distribution.
\end{theorem}
\begin{proof}
  The result follows directly from the classical central limit theorem.
\end{proof}

\chapter{Citing References}

ThuThesis supports both BibTeX and BibLaTeX. This example uses BibTeX with the
\texttt{natbib} package.

\section{Numeric Citations}

Numeric citations are the default for graduate theses, and reference numbers are
normally printed as superscripts. Ordinary citations, author names in prose,
citations with page numbers, and multiple citations can be written with
\verb|\cite|, \verb|\citet|, and \verb|\citep|, for example
\cite{zhangkun1994,zhukezhen1973}.

\section{Author--Year Citations}

Author--year citations can be selected in the document configuration without
changing the bibliography database. Every item in the reference list should
normally be cited in the text. This example includes the complete supplied
database so that the bibliography layout can be inspected.
\nocite{*}

\bibliography{ref/refs}

\appendix
\chapter{Supplementary Material}

An appendix contains material closely related to the thesis that would disrupt
the organization or logic of the main text if included there, such as important
data tables, programs, or statistical tables. It supplements the main body and
may be included as needed.

Figures, tables, and mathematical expressions in an appendix are numbered
separately from those in the main text and carry the appendix identifier.

\section{Figures}
\begin{figure}
  \centering
  
\includegraphics[width=0.6\linewidth]{example-image-a.pdf}
  \caption{Example figure in an appendix}
\end{figure}

\section{Tables}
\begin{table}
  \centering
  \caption{Example table in an appendix}
  \begin{tabular}{ll}
    \toprule
    File name & Description \\
    \midrule
    thuthesis.dtx   & Template source, including documentation and comments \\
    thuthesis.cls   & Template class file \\
    thuthesis-*.bst & BibTeX bibliography style files \\
    thuthesis-*.bbx & BibLaTeX bibliography style files \\
    thuthesis-*.cbx & BibLaTeX citation style files \\
    \bottomrule
  \end{tabular}
\end{table}

\section{Mathematical Expressions}
\begin{equation}
  \frac{1}{2\uppi\symup{i}}\int_\gamma f
  =\sum_{k=1}^m n(\gamma;a_k)\mathscr{R}(f;a_k).
\end{equation}

\section{Citations}
An appendix citation such as \cite{dupont1974bone} can be accompanied by another
reference such as \cite{zhengkaiqing1987}.

\backmatter

\begin{acknowledgements}
  I sincerely thank my supervisor, Professor ×××, and Associate Professor ×× of
  the Department of Physics for their attentive guidance. Their teaching by word
  and example will benefit me throughout my life.

  During nine months of collaborative research in the Department of Chemistry at
  the Massachusetts Institute of Technology, I received generous guidance and
  assistance from Professor Robert Field, for which I am deeply grateful.

  I thank Professor ×××, director of the ××××× Laboratory, as well as all the
  teachers and fellow students in the laboratory for their warm help
  and support.

  This research was supported by the National Natural Science Foundation of
  China, which is gratefully acknowledged.
\end{acknowledgements}

\statement

\begin{resume}
  \section*{Curriculum Vitae}
  Born in a county in Sichuan Province on Month ×, 197×.
  \section*{Academic Work Completed During the Degree Program}
  \subsection*{Research Papers}
  \begin{achievements}
  \item Yang Y, Ren T L, Zhang L T, et al. Miniature microphone with
    silicon-based ferroelectric thin films[J]. Integrated
    Ferroelectrics, 2003, 52:229--235.
  \end{achievements}
  \subsection*{Patents}
  \begin{achievements}
  \item Ren T L, Yang Y, Zhu Y P, et al. Method for controlling
    domain-polarized regions and connecting electrodes in a
    silicon-based ferroelectric microacoustic sensor[P].
  \end{achievements}
\end{resume}

\begin{comments}[name={Comments from the Thesis Supervisor}]
  The thesis presents~\ldots
\end{comments}

\begin{resolution}
  The thesis presents~\ldots

  Its principal original contributions are as follows:

  1.~\ldots

  2.~\ldots

  3.~\ldots

  The work demonstrates that the author possesses knowledge of ××××× and the
  ability to ××××. The thesis ××××, and the defense ××××.

  By a vote of × to × (or unanimously), the defense committee approved the
  thesis defense and recommends that ××× (name) be awarded the Doctor's/Master's
  degree in ××× (discipline category).
\end{resolution}

\end{document}
